\section{One Factor hybrid Hull-White/Dupire model}
Assume now that $a_t \equiv 1$. We consider the particular case where $r_t$ follows the Hull-White model:

\begin{equation}\label{eq: hull_white}
dr_t = \kappa\big(\theta(t)-r_t\big)dt+\sigma_r dB_t
\end{equation}
with $B$ is $\mathbb{Q}$-Brownian motion, correlated with $W$, the correlation being $\rho$.

The function $ \theta(t)$ is chosen such that the zero coupon curve observed at $t = 0$ is matched:
$$\theta(t) = \frac{1}{\kappa}\frac{\partial f^{mkt}}{\partial t}(0,t)+f^{mkt}(0,t)+\frac{\sigma_r^2}{2\kappa^2}(1-e^{-2\kappa t})  $$
where
$$f^{mkt}(0,t)=-\frac{\partial\ln P_{0t}^{mkt}}{\partial t} = r_t^0$$

The bond prices are given by:
\begin{equation}\label{eq: bondPrice_hullWhite}
P_{tT}= \frac{P_{0T}^{mkt}}{P_{0t}^{mkt}}\exp\bigg( -A(t,T)(r_t-r_t^0)-\frac{\sigma_r^2}{4\kappa^3}(e^{-\kappa T}-e^{-\kappa t})^2(e^{2\kappa t}-1)     \bigg)
\end{equation}
where
$$A(t,T)= \frac{1-e^{ -\kappa(T-t) }}{\kappa}$$
and the volatility of the zero-coupon bond is 
$$\sigma_P^T(t)=-\sigma_r A(t,T)$$\\

In view of simulating all processes under the terminal forward measure $\mathbb{Q}^T$ ,
we aim at writing $r_t -r_t^0$ in terms of the $\mathbb{Q}^T$-Brownian motion
$$ B^T_t=B_t-\int_0^t\sigma_P^T(s)ds=B_t+\sigma_r\int_0^tA(s,T)ds $$

After computations the explicit solution to (\ref{eq: hull_white}) under the $\mathbb{Q}^T$-measure is:
\begin{equation}\label{eq: hullWhiteSolution}
r_t -r_t^0= e^{-\kappa t}\bigg(\sigma_r\int_0^t e^{\kappa s}dB_s^T-\frac{\sigma_r^2}{\kappa^2}\sinh(\kappa t)\big(1-e^{-\kappa(T-t)} \big)   \bigg)
\end{equation}

\subsection{Malliavin representation under Hull-White}\label{subsec: malliavin1F}
Under these assumptions, (\ref{eq: malliavinRepresentation}) reads:

\begin{equation}\label{eq: malliavinHullWhite}
\frac{ \mathbb{E}^{\mathbb{Q}}[ D_{0t}(r_t-r^0_t)\mathbb{I}_{S_t>K} ]}{ \frac{1}{2}K\partial^2_KC(t,K) } = 
2\sigma_re^{-\kappa t}\mathbb{E}^{\mathbb{Q}^t}[ V_t(\rho U_t + \Theta_t\Xi_t-\Lambda_t)|S_t = K]
\end{equation}

where:
\begin{align*}
  \frac{dV_t}{V_t} &= S_t\partial_S\sigma(t,S_t)\{dW_t^T + [\rho\sigma_P^T(t,r_t)- \sigma(t,S_t) ]dt\} ,\ \ V_0=1\\
  U_t &= \int_0^te^{\kappa s}\frac{\sigma(s,S_s)}{V_s}ds\\ \numberthis \label{eq: malliavinWeights}
  \Theta_t &= \int_0^t\frac{e^{-\kappa s}}{V_s}ds \\  
  \Lambda_t &= \int_0^t\sigma_re^{2\kappa s}\Theta_s ds\\
  \Xi_t &= \sigma_r\frac{e^{2\kappa t}-1}{2\kappa}
\end{align*}

The computation of (\ref{eq: malliavinHullWhite}) for all $t$, requires the simulation of only 3 processes
$f_t$, $r_t$, $V_t$, and 3 integrals $U_t$, $\Theta_t$, $\Lambda_t$.
The local volatility (\ref{eq: localVolatility_forward}) can be now rewritten as:

\begin{align*}
\sigma(t,K,\mathbb{Q}_t^T)^2=  \sigma_{Dup}(t,K)^2 &-2\sigma_r\rho e^{-\kappa t}\frac{\mathbb{E}^{\mathbb{Q}^T}[P_{tT}^{-1}V_tU_t|S_t=K]}{\mathbb{E}^{\mathbb{Q}^T}[P_{tT}^{-1}|S_t=K]} \\ 
 &  -2\sigma_r e^{-\kappa t}\frac{\mathbb{E}^{\mathbb{Q}^T}[P_{tT}^{-1}V_t(\Theta_t\Xi_t-\Lambda_t)|S_t=K]}{\mathbb{E}^{\mathbb{Q}^T}[P_{tT}^{-1}|S_t=K]} \numberthis \label{eq: localVolatility_forward_HullWhite}
\end{align*}

\subsection{The Particle Method - Hull White }
In the particle algorithm for the hybrid Hull-White/Dupire model, a particle can then be described by $(f_t, r_t- r^0_t )$. Simulate the $T$-forward price $f_t = S_t/P_{tT}$ under the final forward measure as follows:
$$ df_t^i=  f_t^i\sigma_N(t,f_t^iP_{tT}^i)dW_t^i+f_t^i\sigma_r A(t,T)dB_t^i $$
where $B^i$ and $W^i$ are $\mathbb{Q}^T$-Brownian motions.\\
Using the representation (\ref{eq: localVolatility_forward_HullWhite}) for the local volatility, we define:

\begin{align*}
\sigma_N(t,K)^2 = \sigma_{Dup}(t,K)^2 & -2\sigma_r\rho e^{-\kappa t}\frac{\sum_{i=1}^N \Pi_t^{i} V_t^iU_t^i \delta_{t,N}(S_t^i-K) }{\sum_{i=1}^N \Pi_t^{i} \delta_{t,N}(S_t^i-K)} \\
&  -2\sigma_r e^{-\kappa t}\frac{\sum_{i=1}^N \Pi_t^{i}V_t^i(\Theta^i_t\Xi^i_t-\Lambda^i_t)\delta_{t,N}(S_t^i-K) }{\sum_{i=1}^N \Pi_t^{i} \delta_{t,N}(S_t^i-K)} \numberthis \label{eq: discrete_localVolatility_forward_HullWhite}
\end{align*}

with:

$$ \Pi_t^{i} = \exp\bigg( A(t,T)(r_t^i-r_t^0)+\frac{\sigma_r^2}{4\kappa^3}(e^{-\kappa T}-e^{-\kappa t})^2(e^{2\kappa t}-1)     \bigg) $$

\subsection{Time discretization}
As anticipated before, at each iteration of the algorithm and for each particle, we need to simulate 3 processes
$f_t$, $r_t$, $V_t$, and compute 3 integrals $U_t$, $\Theta_t$, $\Lambda_t$.

\subsubsection{Processes: $f_t$ and $r_t$}
Let $\{t_j\}$ be a discretization of $[0,T]$ and consider a single particle.\\
In order to simulate $(f_{t_{j+1}}, r_{t_{j+1}} - r^0_{t_{j+1}})$ from $(f_{t_{j}}, r_{t_{j}} - r^0_{t_{j}})$, one can observe that, from (\ref{eq: hullWhiteSolution}):

\begin{align*}
r_{t_{j+1}}& - r^0_{t_{j+1}}  =  e^{-\kappa(t_{j+1}-t_j) }(r_{t_{j}} - r^0_{t_{j}})+\sigma_r\int_{t_j}^{t_{j+1}}e^{-\kappa(t_{j+1}-t) }dB_t^T \\
& - \frac{\sigma^2_r}{\kappa^2}e^{-\kappa t_{j+1} }\big( \sinh(\kappa t_{j+1})(1-e^{-\kappa(T-t_{j+1}) })-\sinh(\kappa t_{j})(1-e^{-\kappa(T-t_{j}) }) \big) \numberthis \label{eq: discretization_interestRate}
\end{align*}

and use the following approximation scheme for $f_t$:

\begin{align*}
f_{t_{j+1}}& = f_{t_{j}} \exp\bigg(  \sigma(t_j,P_{t_j T}f_{t_j})(W_{t_{j+1}}^T-W^T_{t_{j}})+\sigma_r\int_{t_j}^{t_{j+1}} \frac{1-e^{-\kappa(T-t)}}{\kappa}dB_t^T \\
& -\frac{1}{2}(t_{j+1}-t_{j})\big(\sigma(t_j,P_{t_j T}f_{t_j})^2+2\rho\sigma_r\sigma(t_j,P_{t_j T}f_{t_j})I^1_j + \sigma_r^2I_j^2\big)  \    \bigg) \numberthis \label{eq: discretization_forwardPrice}
\end{align*}
with 
\begin{align*}
& I^1_j = \frac{\kappa(t_{j+1}-t_{j})-e^{-\kappa T }(e^{\kappa t_{j+1} }-e^{\kappa t_{j} }) }{\kappa^2 (t_{j+1}-t_{j})} \\
& I^2_j = \frac{\kappa(t_{j+1}-t_{j})  -2e^{-\kappa T }(e^{\kappa t_{j+1} }-e^{\kappa t_{j} }) +\frac{1}{2}e^{-2\kappa T }(e^{2\kappa t_{j+1} }-e^{2\kappa t_{j} }) }{\kappa^3(t_{j+1}-t_{j})}
\end{align*}
Then the simulation of $(f_{t_{j+1}}, r_{t_{j+1}} - r^0_{t_{j+1}})$ boils down to simulate exactly the Gaussian vector
\[\bigg( W_{t_{j+1}}^T-W^T_{t_{j}},  \int_{t_j}^{t_{j+1}}e^{-\kappa(t_{j+1}-t) }dB_t^T  ,\int_{t_j}^{t_{j+1}} \frac{1-e^{-\kappa(T-t)}}{\kappa}dB_t^T \bigg)  \]
whose covariance matrix is $(t_{j+1}-t_{j})\Sigma_j$ where
\[\Sigma_j= \left( \begin{array}{ccc}
1 & \rho J_j^1 & \rho I_j^1 \\ 
\rho J_j^1 & J_j^2 & L_j \\ 
\rho I_j^1 & L_j & I^2_j \end{array} \right)\] 
with
\begin{align*}
& J_j^1 = \frac{1-e^{-\kappa(t_{j+1}-t_j) }}{\kappa(t_{j+1}-t_j)} \\
& J_j^2 = \frac{1-e^{-2\kappa(t_{j+1}-t_j) } }{2\kappa(t_{j+1}-t_j)} \\
& L_j = \frac{J_j^1}{\kappa}-e^{-\kappa(T + t_{j+1}) } \frac{ e^{2\kappa t_{j+1} }-e^{2\kappa t_{j} }  }{2\kappa^2(t_{j+1}-t_{j})}
\end{align*}
\subsubsection{Process: $V_t$}
For $V_t$ we shall start with a plain Euler discretization scheme. Refinements might follow.\\
The computation of the partial derivative $\partial_S\sigma$ in $V_t$ will be addressed in dependence on the chosen interpolation technique.
\subsubsection{Integrals: $U_t$, $\Theta_t$, $\Lambda_t$}
First note that $V_t$ is computed over the time grid $\{t_j\}$, hence it is a discrete vector of values.
Therefore we set $ V_t = (V_{t_{j+1}}+V_{t_{j}})/2 $ for all $t\in [t_j,t_{j+1}]$. Same reasoning applies for $\Theta_t$.\\
From (\ref{eq: malliavinWeights}) we have: 

\begin{align*}
\Theta_{t_{j+1}} & =  \int_0^{t_{j+1}}\frac{e^{-\kappa s}}{V_s}ds = \Theta_{t_{j}} + \int_{t_j}^{t_{j+1}}\frac{e^{-\kappa s}}{V_s}ds = 
\Theta_{t_{j}} + \frac{2}{(V_{t_{j+1}}+V_{t_{j}})}\int_{t_j}^{t_{j+1}}e^{-\kappa s}ds = \\
 & = \Theta_{t_{j}} - \frac{2}{(V_{t_{j+1}}+V_{t_{j}})}\frac{1}{\kappa}(e^{-\kappa t_{j+1}}-e^{-\kappa t_{j}} )\\ \\
\Lambda_{t_{j+1}} &= \int_0^{t_{j+1}}\sigma_re^{2\kappa s}\Theta_s ds = \Lambda_{t_{j}} + \int_{t_j}^{t_{j+1}}\sigma_re^{2\kappa s}\Theta_s ds = \Lambda_{t_{j}} + \frac{(\Theta_{t_{j+1}} +\Theta_{t_{j}} )}{2}\sigma_r\int_{t_j}^{t_{j+1}}e^{2\kappa s} =\\
& = \Lambda_{t_{j}} + \frac{(\Theta_{t_{j+1}} +\Theta_{t_{j}} )}{4\kappa}\sigma_r(e^{2\kappa t_{j+1}} - e^{2\kappa t_{j}})
\end{align*}
\\
A more careful analysis is needed for $U_t$ as it involves the integration of the local volatility function.
From (\ref{eq: malliavinWeights}) we have:
$$  U_{t_{j+1}} = \int_0^{t_{j+1}}e^{\kappa s}\frac{\sigma(s,S_s)}{V_s}ds =  U_{t_{j}} +  \int_{t_{j}}^{t_{j+1}}e^{\kappa s}\frac{\sigma(s,S_s)}{V_s}ds$$
From step 4 in the algorithm at section \ref{subsec: algorithm}, we have that $\sigma(t,K)\equiv \sigma_N(t_j,K)$ for all $t\in [t_j,t_{j+1}]$, hence we can rewrite the integral as: 
$$ U_{t_{j+1}} = U_{t_{j}} +  \frac{2}{(V_{t_{j+1}}+V_{t_{j}})}\int_{t_{j}}^{t_{j+1}}e^{\kappa s}\sigma(t_j,S_s)ds $$

\IfFileExists{volatility.jpg}{\centerline{\includegraphics[scale=0.5]{volatility.jpg}}}{\centerline{\fbox{\parbox{0.6\textwidth}{\centering\small [figure \texttt{volatility.jpg}: place the file next to this source]}}}}
\bigskip
Now we have two possibilities:
\begin{itemize}
\item approximate $\sigma(t_j,S_s)$ as $[ \sigma(t_j,S_{t_{j+1}}) + \sigma(t_j,S_{t_j}) ]/2$ or $\sigma(t_j,\frac{S_{t_{j+1}} +S_{t_{j}}}{2}) $ and analytically solve the remaining integral
\item apply numerical integration techniques to the integral 
\end{itemize}
The advantage of the first approach is that it is faster to compute and does not depend on the interpolation of the local volatility but it could imply non-negligible errors in the case $S_{t_j}$ is far from $S_{t_{j+1}}$.
The numerical integration might be slower to compute and one should pay attention at the interpolation algorithm used for the local volatility. However it provides a more accurate approximation of the integral, hence we may choose it above the former.

\subsubsection*{Trapezoidal rule for $U_t$}
If using a trapezoidal rule, $U_{t_{j+1}}$ reads:
$$ U_{t_{j+1}} = U_{t_{j}} +  \frac{(t_{j+1}-t_j)}{(V_{t_{j+1}}+V_{t_{j}})}\big[\sigma(t_j,S_{t_j})e^{\kappa t_j}+ \sigma(t_j,S_{t_{j+1}})e^{\kappa t_{j+1}} \big] $$

\subsubsection*{Cavalieri-Simpson rule for $U_t$}
If using a Cavalieri-Simpson rule, $U_{t_{j+1}}$ reads:
$$ U_{t_{j+1}} = U_{t_{j}} +  \frac{1}{3}\frac{(t_{j+1}-t_j)}{(V_{t_{j+1}}+V_{t_{j}})}\big[\sigma(t_j,S_{t_j})e^{\kappa t_j}+  4\sigma\big(t_j,\frac{S_{t_{j+1}} +S_{t_{j}}}{2}\big)e^{\kappa \frac{t_{j+1}+t_{j}}{2}}   +\sigma(t_j,S_{t_{j+1}})e^{\kappa t_{j+1}} \big] $$
\smallskip

\subsection{The Algorithm}\label{subsec: algorithm}
Let $\{t_j\}$ be a discretization of $[0,T]$:

\begin{itemize}
\item[1.] Initialize $j=1$ and set $\sigma_N(t,K)=\sigma_{Dup}(0,K)$ on the first time step
\item[2.] Simulate the $N$ particles  $\{(f_t^i, r_t^i- r^0_t) \}_{1\leq i\leq N}$ from  $t_{j}$ to  $t_{j+1}$ using  (\ref{eq: discretization_interestRate}) and (\ref{eq: discretization_forwardPrice}) 
\item[3.] For each particle, simulate $V_{t_{j+1}}^i$ and compute the integrals $U_{t_{j+1}}^i$, $\Theta_{t_{j+1}}^i$, $\Lambda_{t_{j+1}}^i$ according to (\ref{eq: malliavinWeights})
\item[4.] Sort the particles according to the spot value. For $K\in G_{K,t_j}$, find the smallest index $i(K)$ and the largest $j(K)$ for which $\delta_{t,N}(S_t^i-K) > \eta$. Compute the local volatility.
\begin{itemize}
\item Interpolate it and extrapolate it
\item Set $\sigma(t,K)\equiv \sigma_N(t_j,K)$ for all $t\in [t_j,t_{j+1}]$
\end{itemize}
\item[5.] Set $j:=j+1$, iterate steps 2,3 and 4 up to maturity T
\end{itemize}

\subsection{An equivalent one-process form}\label{subsec: oneProcess}
The five objects $V_t,U_t,\Theta_t,\Lambda_t,\Xi_t$ can be replaced by a single process. Define
$$ \ell_t := \sigma_r e^{-\kappa t}\,V_t\big(\rho U_t+\Theta_t\Xi_t-\Lambda_t\big). $$
Since $\rho U_t+\Theta_t\Xi_t-\Lambda_t$ has finite variation, with
$d(\rho U_t+\Theta_t\Xi_t-\Lambda_t)=V_t^{-1}\big(\rho e^{\kappa t}\sigma(t,S_t)+e^{-\kappa t}\Xi_t\big)dt$,
It\^o's formula gives
\begin{equation}\label{eq: oneProcess}
d\ell_t=\ell_t\Big[g(t,S_t)\big(dW_t-\sigma(t,S_t)dt\big)-\kappa\,dt\Big]
+\Big[v(t)+\rho\,\sigma_r\,\sigma(t,S_t)\Big]dt,\qquad \ell_0=0,
\end{equation}
where
$$ g(t,S)=S\,\partial_S\sigma(t,S),\qquad v(t)=\sigma_r^2\,\frac{1-e^{-2\kappa t}}{2\kappa}, $$
and $W$ is the $\mathbb{Q}$-Brownian motion, $dW_t=dW^T_t+\rho\,\sigma_P^T(t)dt$. Note that $v(t)$ is the variance of the Hull-White state at time $t$.
The right hand side of (\ref{eq: malliavinHullWhite}) is then $2\,\mathbb{E}^{\mathbb{Q}^t}[\ell_t\,|\,S_t=K]$ and (\ref{eq: discrete_localVolatility_forward_HullWhite}) becomes
\begin{equation}\label{eq: oneProcessParticle}
\sigma_N(t,K)^2=\sigma_{Dup}(t,K)^2-2\,\frac{\sum_{i=1}^N\Pi^i_t\,\ell^i_t\,\delta_{t,N}(S^i_t-K)}{\sum_{i=1}^N\Pi^i_t\,\delta_{t,N}(S^i_t-K)} .
\end{equation}
One Euler step of (\ref{eq: oneProcess}) per particle replaces the simulation of $V_t$ and the three quadratures of the previous section. This is the form that extends to the three-factor model of Section \ref{sec: threeFactor}: there, $\ell_t$ is the logarithmic sensitivity of the spot to a perturbation of the rate driver along the Hull-White kernel.
